\documentclass[answers,12pt,addpoints]{exam} 
\usepackage{import}

\import{C:/Users/prana/OneDrive/Desktop/MathNotes}{style.tex}

% Header
\newcommand{\name}{Pranav Tikkawar}
\newcommand{\course}{01:XXX:XXX}
\newcommand{\assignment}{Homework n}
\author{\name}
\title{\course \ - \assignment}

\begin{document}
\maketitle

\newpage

Base model structure:

1-D Spatial, Periodic Boundary on interval $[0,L]$\\
Matern Spatial Covariance\\
Can be written as $X_t(s) = \sum_{j=1}^\infty A_j(t) \cos\left(\frac{2\pi j s}{L}\right) + B_j(t) \sin\left(\frac{2\pi j s}{L}\right)$\\
Temporally Markov: $Z_{t+1} = M Z_t + \epsilon_t$\\
where $\epsilon_t \sim \mathcal{N}(0, \Sigma_\epsilon)$ and $Z_t = (A_1(t), B_1(t), A_2(t), B_2(t), \ldots)^T$\\
$\Sigma_\epsilon$ is correlated across the different frequencies, but not across time.
We want time covariance to be like $\Cov(A_{j,t}, A_{j,t+k}) = \sigma_j^2 \rho_j^k$ and $\Cov(B_{j,t}, B_{j,t+k}) = \sigma_j^2 \rho_j^k$ where $\rho_j$ is a non-constant function of $j$ that decays faster for higher $j$ to control the smoothness of the process in space. Probably want $\rho_j = \exp(-\lambda j)$ for some $\lambda > 0$.


\begin{questions}
\question How does the matern covariance function look like on a perodic domain: Do I just sum over the period and get a new covariance function?

\question Are there some restrictions of the parameters of the matern covariance function to ensure has nice form on the periodic domain?

\question Does this model keep the same smoothness properties as the original matern covariance function, probably not, but how does it change the smoothness properties?

\question What do I expect to see in the high and low frequency components of the covariance function? 
I am assuming that the high frequency components will decay faster than the low frequency components. 
\end{questions}


We will need to observe it everywhere, so then it would make sense to be markov. 

Think about how I would simulate this model. Literally just simulate fourier coefficients. 

Choice of $\Sigma_i^2$ for some algebracitc proety and smoothnes

Also think about hwo to get, some form of drift, and look at cross correlaiton. Thnk of a bivariate AR(1) process for the fourier coefficients, and then see how that translates to the spatio-temporal process. Also see how this translates to phase and amplitude of the fourier coefficients.

\end{document}